\begin{abstract}
Synthetic image detectors report near-perfect benchmark scores but fail when applied outside their training conditions.
Prior work attributes this to distribution shift in generative models, but we argue a simpler explanation has been overlooked: popular benchmarks encode the real-versus-fake label in file format and resolution, and detectors learn the shortcut instead of generative signatures.
We propose evaluating detectors across benchmarks that vary in shortcut severity, and introduce format normalization as a control to separate metadata shortcuts from pixel-level ones.

We test six public detectors on CIFAKE (no format or resolution bias), Defactify (resolution bias only), and GenImage (severe format and resolution bias).
On CIFAKE, all six detectors score below chance, with several calling real images fake more often than fakes.
On GenImage, which shares the JPEG$=$real, PNG$=$fake pattern from their training data, AUC jumps by an average of $0.53$ across five detectors.
Re-encoding every image to $256\times256$ PNG does not recover CIFAKE performance, which means the shortcut lives in pixel statistics, not metadata.
One detector, B-Free, behaves differently: its DINOv2 features rise from $0.497$ to $0.637$ on CIFAKE, suggesting it picks up visual quality rather than compression artifacts.
Below-chance performance on an artifact-free benchmark is a clean diagnostic for shortcut-dependent detectors.

To probe causality, we conduct a format-swap experiment: re-encoding real images as PNG and fake images as JPEG reverses the format association.
FreqNet drops from AUC $0.977$ to $0.503$ under this swap; NPR inverts to $0.381$.
Frequency band ablation further localizes the shortcut: zeroing the low-frequency band ($r < 0.2\,r_{\max}$) degrades all six detectors substantially on the biased benchmark, with NPR falling by $0.56$ AUC and UnivFD by $0.40$.
\end{abstract}
